% Build with: make
% LuaLaTeX; no figures, bibliography database or research data are required.
\documentclass[12pt,a4paper]{article}
\usepackage[margin=21mm,top=18mm,bottom=19mm]{geometry}
\usepackage{fontspec}
\setmainfont{Liberation Serif}
\usepackage{babel}
\babelprovide[main,import]{russian}
\babelprovide[import]{english}
\usepackage{microtype}
\usepackage{array,tabularx,booktabs}
\usepackage{titlesec}
\usepackage[unicode,colorlinks=true,urlcolor=blue,linkcolor=black]{hyperref}
\hypersetup{
  pdftitle={Отчёт о научно-исследовательской работе с 20 мая по 20 ноября 2026 года},
  pdfsubject={Развитие GRTeclyn и исследования скалярных полей в общей теории относительности},
  pdflang={ru-RU}
}
\urlstyle{same}
\titleformat{\section}{\normalsize\bfseries}{}{0pt}{}
\titlespacing*{\section}{0pt}{9pt}{4pt}
\setlength{\parindent}{0pt}
\setlength{\parskip}{5pt}
\setlength{\emergencystretch}{2em}
\linespread{1.025}
\newcolumntype{L}[1]{>{\raggedright\arraybackslash}p{#1}}
\newcolumntype{Y}{>{\raggedright\arraybackslash}X}
\setlength{\tabcolsep}{5pt}
\renewcommand{\arraystretch}{1.08}
\pagestyle{plain}
\raggedbottom
\clubpenalty=10000
\widowpenalty=10000

\begin{document}

\begin{center}
  {\large\bfseries Отчёт о научно-исследовательской работе\par}
  \vspace{3pt}
  {с 20 мая по 20 ноября 2026 года\par}
\end{center}
\vspace{-3pt}

Созданы стационарный решатель и автоматический пайплайн расчётов на GPU. Они позволяют находить решения уравнений Эйнштейна со скалярной материей и проверять, как эти решения меняются со временем.

\section*{Стационарный решатель и автоматический пайплайн расчётов}

Создан стационарный решатель, который совместно определяет распределение скалярного поля и геометрию пространства-времени. С его помощью построены стационарные вращающиеся кротовые норы для расчётов их эволюции. Затем добавлены модели бозонных звёзд и скалярных колец с разными потенциалами поля. Решатель оформлен как самостоятельный программный пакет.

Создан автоматический пайплайн расчётов: решение стационарной задачи, подготовка начальных данных, запуск трёхмерной эволюции в GRTeclyn на GPU и анализ результатов. Система сохраняет параметры и версии кода, следит за ходом расчёта, строит графики и анимации. Это позволяет проводить серии исследований без ручного переноса данных между программами.

Проверки встроены в обе части системы. Стационарные решения подставляются в дополнительные уравнения Эйнштейна, которые не использовались при их построении. При эволюции контролируются гамильтонова и импульсные связи. Расчёты повторяются при повышении разрешения: ошибки должны уменьшаться, а физические результаты --- сходиться.

Авторы исходной работы о вращающихся кротовых норах подтвердили найденную нами опечатку в уравнении, отметили отличное согласие расчётов и поблагодарили за проверку. Их численные результаты не затронуты: в коде использовалось верное уравнение.

В основной репозиторий GRTeclyn приняты наши исправления расчёта связей и извлечения гравитационных волн при наличии материи. Они устраняют использование незаполненных ячеек на границах вычислительных блоков: \href{https://github.com/GRTLCollaboration/GRTeclyn/pull/241}{GRTLCollaboration/GRTeclyn, PR~\#241}.

Код, данные и программы анализа: \url{https://github.com/Nikchik-coder/GRTeclyn}.

\section*{Основные научные результаты}

\textit{Экзотическая материя и диполь Бонди.} Реализована модель скалярной материи, которая входит в уравнения Эйнштейна с отрицательным знаком. Построены бозонные звёзды с отрицательной гравитационной массой, сохраняющие структуру при длительной эволюции. По известной нам литературе, впервые рассчитана полная трёхмерная эволюция таких звёзд и диполя Бонди: пара тел с положительной и отрицательной массами ускоряется как единое целое при почти постоянном расстоянии между ними. Результат проверен контрольными расчётами и исследованием сходимости.

\textit{Вращающиеся кротовые норы.} Впервые исследована их нелинейная трёхмерная эволюция. Вращение замедляет неустойчивость. У наиболее быстро вращающейся модели рост радиальных возмущений не обнаружен на интервале около 28 характерных времён роста неустойчивости невращающейся норы. Этот результат получен при трёх разрешениях для осесимметричных начальных данных. Показано также, что коллапс вращающейся норы излучает гравитационные волны без специально добавленного несферического возмущения.

\newpage

\section*{Статьи на arXiv}

По результатам работы опубликованы три препринта Н.~М.~Широкова и И.~Начевского.

{\small
\begin{tabularx}{\textwidth}{@{}L{36mm}YL{34mm}@{}}
\toprule
\textbf{Работа и ссылка} & \textbf{Результат} & \textbf{Журнал и статус} \\
\midrule
Диполь Бонди\par
\url{https://arxiv.org/abs/2608.24577}
& Эволюция звёзд с отрицательной гравитационной массой и самоускорение пары с массами разных знаков в полной общей теории относительности.
& \textit{Physics of the Dark Universe}.\par У редактора. \\
\midrule
Вращающиеся кротовые норы\par
\url{https://arxiv.org/abs/2610.29941}
& Замедление неустойчивости при вращении. При быстром вращении горловина сохраняет размер за время расчёта. Коллапс сопровождается гравитационными волнами.
& \textit{European Physical Journal C}.\par Статья подана. \\
\midrule
Двойные кротовые норы\par
\url{https://arxiv.org/abs/2610.13206}
& Столкновения и пролёты пар кротовых нор, образование общего горизонта, гравитационное и скалярное излучение.
& \textit{General Relativity and Gravitation}.\par Статья подана. \\
\bottomrule
\end{tabularx}
}

\section*{Другие исследования и дальнейшая работа}

Для оптимизации параметров применены методы эволюционного поиска, используемые в машинном обучении. Найденные решения проверяются теми же уравнениями и расчётами эволюции.

{\small
\begin{tabularx}{\textwidth}{@{}L{31mm}YL{31mm}@{}}
\toprule
\textbf{Направление} & \textbf{Что сделано} & \textbf{Состояние работы} \\
\midrule
Минимизация экзотической материи
& Для выбранной вращающейся норы подбор потенциала скалярного поля уменьшил количество экзотической материи примерно на 61\% при тех же размере горловины и угловом моменте. Оптимизированные решения проверены в трёхмерной эволюции на GPU.
& Основные расчёты завершены.\par Статья подготовлена. \\
\midrule
Скалярные кольца
& Построены вращающиеся кольца с эргообластями и пары колец с противоположным вращением. Исследовано движение частиц в их поле. Для отдельного кольца без эргообласти выполнена эволюция на GPU; начальный этап проверен на трёх разрешениях.
& Готовится статья.\par Продолжается исследование устойчивости. \\
\midrule
Кротовая нора с кольцом
& Получены совместные решения для кротовой норы и вращающегося скалярного кольца, в том числе с эргообластью у горловины.
& В работе.\par Поиск моделей с меньшими приливными силами. \\
\midrule
Анализ пространства-времени
& Разработаны метрики и программы для анализа геометрии по данным моделирования. Они оценивают кривизну, изменение геометрии и выигрыш во времени прохождения света по сравнению с плоским пространством (метрика FTL). Расчёт траекторий света в меняющейся геометрии позволяет определять, образуется ли короткий путь и как долго он сохраняется.
& Программы реализованы.\par Проект статьи подготовлен. \\
\bottomrule
\end{tabularx}
}

\end{document}
